%%======================================================================
%%  NT-4b: Literature Extraction & Convention Fixing
%%         for Full Nonlinear Variational Field Equations
%%
%%  This document collects formulas from the literature needed to
%%  derive the full nonlinear field equations delta S_spec/delta g^{mu nu} = 0
%%  on arbitrary curved backgrounds from the SCT spectral action.
%%
%%  RESEARCH DOCUMENT: formulas extracted from cited papers.
%%  No original derivations --- gaps noted explicitly.
%%======================================================================
\documentclass[12pt,a4paper]{article}

%% ---- packages -------------------------------------------------------
\usepackage[utf8]{inputenc}
\usepackage[T1]{fontenc}
\usepackage{lmodern}
\usepackage{amsmath,amssymb,amsthm,mathtools}
\usepackage[margin=2.5cm]{geometry}
\usepackage{booktabs}
\usepackage{enumitem}
\usepackage{hyperref}
\usepackage{xcolor}
\usepackage{fancyhdr}
\usepackage{titlesec}
\usepackage{longtable}
\usepackage{tcolorbox}

%% ---- theorem-like environments --------------------------------------
\newtheorem{proposition}{Proposition}[section]
\newtheorem{lemma}[proposition]{Lemma}
\newtheorem{corollary}[proposition]{Corollary}
\newtheorem{theorem}[proposition]{Theorem}
\theoremstyle{definition}
\newtheorem{definition}[proposition]{Definition}
\newtheorem{convention}[proposition]{Convention}
\theoremstyle{remark}
\newtheorem{remark}[proposition]{Remark}

%% ---- custom commands ------------------------------------------------
\newcommand{\Tr}{\operatorname{Tr}}
\newcommand{\tr}{\operatorname{tr}}
\newcommand{\Rsq}{R_{\mu\nu\rho\sigma}R^{\mu\nu\rho\sigma}}
\newcommand{\Ricsq}{R_{\mu\nu}R^{\mu\nu}}
\newcommand{\Weylsq}{C_{\mu\nu\rho\sigma}C^{\mu\nu\rho\sigma}}
\newcommand{\dd}{\mathrm{d}}
\newcommand{\ii}{\mathrm{i}}
\newcommand{\Id}{\mathrm{Id}}
\newcommand{\erfi}{\operatorname{erfi}}
\DeclareMathOperator{\sgn}{sgn}

%% ---- annotation commands --------------------------------------------
\newcommand{\fromlit}[1]{\textcolor{blue!70!black}{\textsf{[#1]}}}
\newcommand{\oursign}{\textcolor{teal!80!black}{\textsf{[our convention]}}}
\newcommand{\gap}{\textcolor{red!80!black}{\textsf{[GAP]}}}
\newcommand{\needscheck}{\textcolor{orange!80!black}{\textsf{[needs check]}}}
\newcommand{\exactlit}{\textcolor{blue!60!black}{\textsf{[exact from lit.]}}}
\newcommand{\verified}[1]{%
  \par\smallskip\noindent
  \colorbox{green!10}{\parbox{\dimexpr\linewidth-2\fboxsep}{%
    \textcolor{green!50!black}{\textbf{[Verified]}} #1}}%
  \par\smallskip}
\newcommand{\signcorrection}[1]{%
  \par\noindent\colorbox{red!10}{%
    \parbox{\dimexpr\linewidth-2\fboxsep}{%
      \textcolor{red!50!black}{\textbf{[Sign correction]}}
      \textcolor{red!50!black}{#1}%
    }%
  }\par%
}

%% ---- page style -----------------------------------------------------
\pagestyle{fancy}
\fancyhf{}
\fancyhead[L]{\small\textsc{NT-4b: Literature \& Conventions for
Full Nonlinear Field Equations}}
\fancyhead[R]{\small\thepage}
\renewcommand{\headrulewidth}{0.4pt}

%% ---- hyperref -------------------------------------------------------
\hypersetup{
  colorlinks=true,
  linkcolor=blue!60!black,
  citecolor=green!40!black,
  urlcolor=blue!60!black,
  pdftitle={NT-4b: Literature Extraction for Full Nonlinear Field Equations},
  pdfauthor={David Alfyorov},
}

%% ---- title ----------------------------------------------------------
\title{\textbf{NT-4b: Literature Extraction \& Convention Fixing\\[4pt]
       for Full Nonlinear Variational Field Equations}\\[8pt]
       \large From Nonlocal Effective Action to $\frac{\delta S}{\delta g^{\mu\nu}} = 0$
       on Curved Backgrounds}
\author{David Alfyorov}
\date{March 2026 --- Working Document}

%%======================================================================
\begin{document}
\maketitle

\begin{center}
\small\textit{Credits: Aliaksandr Samatyia contributed research-assistance
and workflow support.}
\end{center}

\thispagestyle{empty}

\begin{abstract}
This document collects all formulas, conventions, and computational
recipes from the literature needed to derive the full nonlinear
gravitational field equations from the SCT spectral action
$\Tr(f(D^2/\Lambda^2))$ on arbitrary curved backgrounds.
Unlike the NT-4a linearization (flat background, linear perturbation),
NT-4b requires the variation of a nonlocal curvature-squared action
$\delta S/\delta g^{\mu\nu}$ where the form factors
$F_i(\Box/\Lambda^2)$ do not commute with the metric variation.
The document covers: (1)~the SCT effective action and its conventions
inherited from NT-1 through NT-4a; (2)~the local Stelle field equations
as the limiting benchmark; (3)~the Barvinsky--Vilkovisky formalism for
varying nonlocal actions; (4)~the nonlocal Bach tensor and nonlocal
$R^2$ variation; (5)~key identities for $\delta\Box/\delta g^{\mu\nu}$
and the $\alpha$-insertion formula; (6)~consistency checks (Bianchi
identity, trace, linearized limit, de~Sitter solutions); (7)~a
comprehensive notation table.  All formulas are extracted from the
literature with source annotations; no original derivations are
performed.  Gaps requiring resolution in the derivation phase are noted
explicitly.

\medskip\noindent
\textbf{Key references:}
\begin{itemize}[nosep]
\item Stelle: Phys.\ Rev.\ D \textbf{16} (1977) 953; Gen.\ Rel.\ Grav.\
  \textbf{9} (1978) 353.
\item Barvinsky--Vilkovisky (BV): Phys.\ Rep.\ \textbf{119} (1985) 1;
  Nucl.\ Phys.\ B \textbf{282} (1987) 163;
  B \textbf{333} (1990) 471.
\item Codello--Zanusso (CZ): arXiv:1203.2034 [hep-th].
\item Biswas--Mazumdar--Siegel (BMS): hep-th/0508194.
\item Modesto: arXiv:1107.2403 [hep-th].
\item Calcagni--Modesto--Nardelli (CMN): arXiv:1803.07848;
  also Calcagni--Modesto--Leszczyszyn: arXiv:1803.00561.
\item Buoninfante--Modesto--Shapiro: arXiv:2001.07830.
\item Giacchini--de~Paula~Netto: arXiv:1806.05664 [gr-qc];
  Eur.\ Phys.\ J.\ C \textbf{81} (2021) 741.
\item Deser--Woodard: Phys.\ Rev.\ Lett.\ \textbf{99} (2007) 111301.
\item Avramidi: ``Heat Kernel and Quantum Gravity,'' Springer (2000).
\item Parker--Toms: ``Quantum Field Theory in Curved Spacetime,''
  Cambridge (2009).
\end{itemize}
\end{abstract}

\tableofcontents
\newpage

%%======================================================================
\section{Background: From NT-4a to NT-4b}
\label{sec:background}
%%======================================================================

\subsection{What NT-4a established}

NT-4a derived the \emph{linearized} field equations of the SCT spectral
action on a flat background.  The key results are:
\begin{itemize}[nosep]
  \item \textbf{Modified propagator:}
    $G^{-1}_{\mathrm{TT}} = k^2\,\Pi_{\mathrm{TT}}(z)$ with
    $\Pi_{\mathrm{TT}}(z) = 1 + \frac{13}{60}\,z\,\hat{F}_1(z)$.
  \item \textbf{Scalar mode:}
    $\Pi_s(z,\xi) = 1 + 6(\xi-1/6)^2\,z\,\hat{F}_2(z,\xi)$.
  \item \textbf{Yukawa potential:}
    $V(r)/V_N(r) = 1 - \frac{4}{3}e^{-m_2 r} + \frac{1}{3}e^{-m_0 r}$.
  \item \textbf{Gauge invariance} and \textbf{Bianchi identity}
    verified off-shell.
  \item \textbf{40/40 checks PASS}, 7 Lean~4 identities verified.
\end{itemize}
All these results are \emph{valid only on flat background at linear order}.

\subsection{What NT-4b must achieve}

NT-4b must derive the \textbf{full nonlinear} field equations
\begin{equation}\label{eq:NT4b-goal}
  \frac{1}{\sqrt{g}}\,
  \frac{\delta S_{\mathrm{SCT}}}{\delta g^{\mu\nu}} = 0
\end{equation}
on \emph{arbitrary} curved backgrounds, where $S_{\mathrm{SCT}}$ includes
both the Einstein--Hilbert term and the nonlocal curvature-squared
corrections.

The central difficulty is that the form factors $F_i(\Box/\Lambda^2)$
are functions of the covariant d'Alembertian, which itself depends on
the metric.  Therefore $[\delta/\delta g^{\mu\nu},\, F(\Box)] \neq 0$,
and the variation generates additional ``$\Theta$-type'' terms absent
in the local theory.

\subsection{Gap analysis from NT-4a}

The NT-4a literature document~\cite{NT4aLit} identified the following
gaps relevant to NT-4b:
\begin{enumerate}[label=\textbf{G\arabic*.}]
  \item Full nonlinear equations on curved background
    (Barvinsky--Vilkovisky formalism) --- \textbf{this document}.
  \item FLRW reduction --- deferred to NT-4c.
  \item Covariant gauge-fixing analysis --- needed for NT-4b but
    secondary.
\end{enumerate}


%%======================================================================
\section{The SCT Effective Action: Conventions}
\label{sec:action}
%%======================================================================

\subsection{Sign and curvature conventions}

\begin{convention}[Curvature conventions]
\label{conv:curvature}
\oursign\;
Throughout this project, the conventions are:
\begin{center}
\renewcommand{\arraystretch}{1.3}
\begin{tabular}{ll}
\toprule
\textbf{Quantity} & \textbf{Convention} \\
\midrule
Metric signature & $(-,+,+,+)$ (Lorentzian) or Euclidean $\delta_{\mu\nu}$ \\
Riemann tensor & $R^\rho{}_{\sigma\mu\nu}
  = \partial_\mu\Gamma^\rho_{\nu\sigma}
  - \partial_\nu\Gamma^\rho_{\mu\sigma}
  + \Gamma^\rho_{\mu\lambda}\Gamma^\lambda_{\nu\sigma}
  - \Gamma^\rho_{\nu\lambda}\Gamma^\lambda_{\mu\sigma}$ \\
Ricci tensor & $R_{\mu\nu} = R^\rho{}_{\mu\rho\nu}$ \\
Ricci scalar & $R = g^{\mu\nu}R_{\mu\nu}$ \\
Einstein tensor & $G_{\mu\nu} = R_{\mu\nu} - \frac{1}{2}g_{\mu\nu}R$ \\
d'Alembertian & $\Box = g^{\mu\nu}\nabla_\mu\nabla_\nu
  = \frac{1}{\sqrt{|g|}}\partial_\mu(\sqrt{|g|}\,g^{\mu\nu}\partial_\nu)$
  on scalars \\
Weyl tensor & $C_{\mu\nu\rho\sigma} = R_{\mu\nu\rho\sigma}
  + \frac{2}{d-2}(g_{\nu[\rho}R_{\sigma]\mu} - g_{\mu[\rho}R_{\sigma]\nu})
  + \frac{2}{(d-1)(d-2)}R\,g_{\mu[\rho}g_{\sigma]\nu}$ \\
Gravitational coupling & $\kappa^2 = 8\pi G$ \\
\bottomrule
\end{tabular}
\end{center}
These match the conventions used in
Calcagni--Modesto~\cite{CalcagniModesto2022},
Stelle~\cite{Stelle1977,Stelle1978}, and all preceding
NT-1/NT-1b/NT-2/NT-4a documents.
\end{convention}


\subsection{The full SCT action}

\begin{convention}[Full SCT gravitational action]
\label{conv:full-action}
\oursign\;
The full one-loop SCT gravitational action is:
\begin{equation}\label{eq:full-action}
  S = S_{\mathrm{EH}} + S_{\mathrm{quad}} + S_{\mathrm{matter}},
\end{equation}
where:
\begin{align}
  S_{\mathrm{EH}} &= \frac{1}{2\kappa^2}\int\dd^4 x\,\sqrt{|g|}\,
    (R - 2\Lambda_{\mathrm{cc}}),
  \label{eq:SEH}\\[6pt]
  S_{\mathrm{quad}} &= \frac{1}{16\pi^2}\int\dd^4 x\,\sqrt{|g|}\,
  \Big[
    F_1\!\left(\frac{\Box}{\Lambda^2}\right)\Weylsq
    + F_2\!\left(\frac{\Box}{\Lambda^2},\xi\right) R^2
  \Big],
  \label{eq:Squad}\\[6pt]
  S_{\mathrm{matter}} &= \int\dd^4 x\,\sqrt{|g|}\,
    \mathcal{L}_{\mathrm{matter}}(g_{\mu\nu}, \psi).
  \label{eq:Smatter}
\end{align}
The cosmological constant term $\Lambda_{\mathrm{cc}}$ comes from the
$a_0$ Seeley--DeWitt coefficient; we keep it general here.
\end{convention}

\begin{remark}[Curvature basis choice]
\label{rem:basis}
The $\{C^2, R^2\}$ basis is preferred over $\{R_{\mu\nu}^2, R^2\}$
because:
\begin{enumerate}[nosep]
  \item $C_{\mu\nu\rho\sigma}$ is traceless, so $C^2$ projects onto
    pure spin-2.
  \item The Gauss--Bonnet identity in $d=4$:
    \begin{equation}\label{eq:GB}
      E_4 = R_{\mu\nu\rho\sigma}R^{\mu\nu\rho\sigma}
      - 4\Ricsq + R^2
    \end{equation}
    is topological ($\int\sqrt{g}\,E_4 = 32\pi^2\chi(M)$) and does
    not contribute to the equations of motion.
  \item Conversion:
    \begin{equation}\label{eq:Weyl-decomp}
      \Weylsq = R_{\mu\nu\rho\sigma}R^{\mu\nu\rho\sigma}
      - 2\Ricsq + \tfrac{1}{3}R^2.
    \end{equation}
\end{enumerate}
\end{remark}

\begin{remark}[Alternative Ricci--scalar basis]
\label{rem:alt-basis}
\fromlit{Calcagni--Modesto 2211.05606, Eq.~(NLGActionGeneral)}
The general nonlocal gravity action in the
$\{R_{\mu\nu\rho\sigma}^2, R_{\mu\nu}^2, R^2\}$ basis reads:
\begin{equation}\label{eq:NLG-general}
  S_{\mathrm{NLG}} = \frac{1}{2\kappa^2}\int\dd^4 x\,\sqrt{|g|}\,
  \big[
    R + R\,\gamma_0(\Box)\,R
    + R_{\mu\nu}\,\gamma_2(\Box)\,R^{\mu\nu}
    + R_{\mu\nu\rho\sigma}\,\gamma_4(\Box)\,R^{\mu\nu\rho\sigma}
  \big].
\end{equation}
Using the Gauss--Bonnet identity~\eqref{eq:GB} to eliminate
$R_{\mu\nu\rho\sigma}^2$, this becomes:
\begin{equation}
  S = \frac{1}{2\kappa^2}\int\dd^4 x\,\sqrt{|g|}\,
  \big[
    R + R\,\tilde{\gamma}_0(\Box)\,R
    + R_{\mu\nu}\,\tilde{\gamma}_2(\Box)\,R^{\mu\nu}
  \big]
  + \text{topological},
\end{equation}
with $\tilde{\gamma}_0 = \gamma_0 - \gamma_4$ and
$\tilde{\gamma}_2 = \gamma_2 + 4\gamma_4$.

To map to our $\{C^2, R^2\}$ basis, one uses
\eqref{eq:Weyl-decomp} to express:
\begin{equation}\label{eq:basis-map}
  \alpha\,\Weylsq + \beta\,R^2
  = \alpha\,R_{\mu\nu\rho\sigma}^2
  - 2\alpha\,\Ricsq
  + \left(\beta + \tfrac{\alpha}{3}\right)R^2.
\end{equation}
Matching: $\gamma_4 = \alpha$, $\gamma_2 = -2\alpha$,
$\gamma_0 = \beta + \alpha/3$.
\end{remark}


\subsection{Input coefficients from Phase~3}

\fromlit{NT-1b Phase~3}
\begin{align}
  \alpha_C &= 16\pi^2 F_1(0) = \frac{13}{120},
  \label{eq:alpha_C}\\
  \alpha_R(\xi) &= 16\pi^2 F_2(0,\xi) = 2\left(\xi-\frac{1}{6}\right)^2.
  \label{eq:alpha_R}
\end{align}
In the $\{R^2, R_{\mu\nu}^2\}$ basis:
\begin{equation}\label{eq:c1c2}
  c_2 = 2\alpha_C = \frac{13}{60}, \qquad
  3c_1 + c_2 = 6\left(\xi-\frac{1}{6}\right)^2.
\end{equation}

\verified{$\alpha_C = 13/120$ and $\alpha_R(\xi) = 2(\xi-1/6)^2$
confirmed by NT-1b Phase~3 (354/354 checks PASS).}

\subsection{The normalized form factors}

\begin{definition}[Normalized form factors]
\label{def:Fhat}
\oursign\;
\begin{equation}
  \hat{F}_i(z) = \frac{F_i(z)}{F_i(0)}, \qquad \hat{F}_i(0) = 1.
\end{equation}
\end{definition}

The master function underlying all form factors is:
\begin{equation}\label{eq:phi-master}
  \varphi(z) = e^{-z/4}\sqrt{\frac{\pi}{z}}\,\erfi\!\left(\frac{\sqrt{z}}{2}\right),
  \qquad \varphi(0) = 1, \quad \varphi'(0) = -\frac{1}{6}.
\end{equation}
The Taylor coefficients are $a_n = (-1)^n n!/(2n+1)!$ (NT-2, verified
entire).


%%======================================================================
\section{Local Benchmark: Stelle's Nonlinear Field Equations}
\label{sec:stelle}
%%======================================================================

The local limit $F_1 \to F_1(0)$, $F_2 \to F_2(0)$ must reproduce
Stelle's fourth-order gravity.  This section documents the exact
nonlinear field equations from Stelle's 1977--1978
papers~\cite{Stelle1977,Stelle1978}.

\subsection{The Stelle action and field equations}

\begin{convention}[Stelle action]
\label{conv:stelle}
\fromlit{Stelle, Phys.\ Rev.\ D \textbf{16} (1977) 953;
Gen.\ Rel.\ Grav.\ \textbf{9} (1978) 353}
\begin{equation}\label{eq:stelle-action}
  S_{\mathrm{Stelle}} = \frac{1}{2\kappa^2}\int\dd^4 x\,\sqrt{|g|}\,
  \left[R + \alpha\,\Weylsq + \beta\,R^2\right].
\end{equation}
\end{convention}

\begin{theorem}[Stelle's nonlinear field equations]
\label{thm:stelle-eom}
\fromlit{Stelle 1978, Eqs.~(3.2)--(3.5)}
The metric variation $\delta S_{\mathrm{Stelle}}/\delta g^{\mu\nu} = 0$
gives:
\begin{equation}\label{eq:stelle-eom}
  \boxed{
    G_{\mu\nu} + 2\alpha\,B_{\mu\nu} + \beta\,H_{\mu\nu}
    = \kappa^2 T_{\mu\nu},
  }
\end{equation}
where the \textbf{Bach tensor} $B_{\mu\nu}$ and the
\textbf{scalar curvature tensor} $H_{\mu\nu}$ are:
\begin{align}
  B_{\mu\nu} &= \nabla^\rho\nabla^\sigma C_{\mu\rho\nu\sigma}
    + \frac{1}{2}R^{\rho\sigma}C_{\mu\rho\nu\sigma},
  \label{eq:bach}\\[6pt]
  H_{\mu\nu} &= 2\nabla_\mu\nabla_\nu R - 2g_{\mu\nu}\Box R
    - \frac{1}{2}g_{\mu\nu}R^2 + 2R\,R_{\mu\nu}.
  \label{eq:H-tensor}
\end{align}
\end{theorem}

\begin{remark}[Alternative forms of the Bach tensor]
\label{rem:bach-alt}
\fromlit{Stelle 1978; Bach 1921}
Using the Bianchi identity, the Bach tensor can be rewritten as:
\begin{equation}\label{eq:bach-alt}
  B_{\mu\nu} = \Box R_{\mu\nu} - \frac{1}{3}\nabla_\mu\nabla_\nu R
    - \frac{1}{6}g_{\mu\nu}\Box R
    + R_{\mu\rho\nu\sigma}R^{\rho\sigma}
    - \frac{1}{3}R\,R_{\mu\nu}
    - \frac{1}{4}g_{\mu\nu}\left(\Ricsq - \frac{1}{3}R^2\right).
\end{equation}
This form makes the fourth-derivative (quasi-linear) structure manifest:
$B_{\mu\nu}$ contains $\Box R_{\mu\nu}$ (four derivatives of $g$) as
its leading term.

The Bach tensor is traceless in $d = 4$:
\begin{equation}\label{eq:bach-trace}
  g^{\mu\nu}B_{\mu\nu} = 0.
\end{equation}
This follows from the conformal invariance of $\int\sqrt{g}\,\Weylsq$.
\end{remark}

\begin{remark}[Properties of $H_{\mu\nu}$]
\label{rem:H-props}
\fromlit{Stelle 1978}
The trace of $H_{\mu\nu}$ is:
\begin{equation}\label{eq:H-trace}
  g^{\mu\nu}H_{\mu\nu} = -6\Box R.
\end{equation}
The divergence satisfies:
\begin{equation}\label{eq:H-div}
  \nabla^\mu H_{\mu\nu} = 0,
\end{equation}
which ensures $\nabla^\mu(G_{\mu\nu} + 2\alpha B_{\mu\nu}
+ \beta H_{\mu\nu}) = 0$ (Bianchi identity for the full equations).
\end{remark}

\signcorrection{Some references define $B_{\mu\nu}$ with the opposite
sign convention.  Our convention follows Stelle (1978) and
Parker--Toms (2009), where $\delta(\Weylsq)/\delta g^{\mu\nu}
= 2B^{\mu\nu}$ (with the factor 2 absorbed into the coefficient
$\alpha$).  Other references (e.g., Oliva--Ray 2010) define
$B_{\mu\nu}$ such that $\delta(\Weylsq)/\delta g^{\mu\nu}
= -2B^{\mu\nu}$.  The sign must be tracked carefully in NT-4b.}


\subsection{Variation of the Gauss--Bonnet term}

\begin{proposition}[Gauss--Bonnet variation in $d=4$]
\label{prop:GB-var}
\exactlit\;
\fromlit{Padmanabhan, ``Gravitation,'' Ch.~15; Lovelock 1971}
In $d = 4$, the Gauss--Bonnet combination $E_4$ is topological:
\begin{equation}\label{eq:GB-var}
  \frac{\delta}{\delta g^{\mu\nu}}\int\dd^4 x\,\sqrt{|g|}\,E_4 = 0.
\end{equation}
This identity is crucial: it means that any basis conversion between
$\{R_{\mu\nu\rho\sigma}^2, \Ricsq, R^2\}$ and $\{\Weylsq, R^2\}$
does not generate additional equations-of-motion contributions.
\end{proposition}


\subsection{Individual curvature-squared variations (local)}

For reference and as building blocks for the nonlocal case, we record
the local metric variations of each curvature-squared invariant.

\begin{proposition}[Metric variation of curvature-squared terms]
\label{prop:curv-sq-var}
\fromlit{Padmanabhan Ch.~15; Parker--Toms App.~B;
Boulware--Deser, Phys.\ Rev.\ Lett.\ \textbf{55} (1985) 2656}
\begin{align}
  \frac{1}{\sqrt{g}}\frac{\delta(\sqrt{g}\,R^2)}{\delta g^{\mu\nu}}
  &= 2\nabla_\mu\nabla_\nu R - 2g_{\mu\nu}\Box R
    + 2R\,R_{\mu\nu} - \frac{1}{2}g_{\mu\nu}R^2
  \notag\\
  &\equiv H_{\mu\nu},
  \label{eq:var-R2}\\[6pt]
  \frac{1}{\sqrt{g}}\frac{\delta(\sqrt{g}\,\Ricsq)}{\delta g^{\mu\nu}}
  &= \Box R_{\mu\nu} + \frac{1}{2}g_{\mu\nu}\Box R
    - \nabla_\mu\nabla_\nu R
    + 2R_{\mu\rho\nu\sigma}R^{\rho\sigma}
    - \frac{1}{2}g_{\mu\nu}\Ricsq,
  \label{eq:var-Ric2}\\[6pt]
  \frac{1}{\sqrt{g}}\frac{\delta(\sqrt{g}\,\Rsq)}{\delta g^{\mu\nu}}
  &= 4\Box R_{\mu\nu} + 4g_{\mu\nu}\Box R
    - 6\nabla_\mu\nabla_\nu R
    + 8R_{\mu\rho\nu\sigma}R^{\rho\sigma}
    - 2g_{\mu\nu}\Ricsq
  \notag\\
  &\quad + \frac{1}{2}g_{\mu\nu}R^2
    - 2R\,R_{\mu\nu}.
  \label{eq:var-Riem2}
\end{align}
\end{proposition}

\begin{remark}[Relation between variations]
Using the Gauss--Bonnet identity~\eqref{eq:GB-var}, the three
variations~\eqref{eq:var-R2}--\eqref{eq:var-Riem2} satisfy:
\begin{equation}
  \frac{\delta(\sqrt{g}\,\Rsq)}{\delta g^{\mu\nu}}
  - 4\frac{\delta(\sqrt{g}\,\Ricsq)}{\delta g^{\mu\nu}}
  + \frac{\delta(\sqrt{g}\,R^2)}{\delta g^{\mu\nu}} = 0.
\end{equation}
This provides a consistency check on all three formulas.
\end{remark}


%%======================================================================
\section{Nonlocal Variation: The Central Problem}
\label{sec:nonlocal-variation}
%%======================================================================

\subsection{Why the variation is nontrivial}

For a local action $S = \int\sqrt{g}\,A\cdot B$ where $A$ and $B$ are
curvature scalars, the variation is simply
$\delta S/\delta g^{\mu\nu} = (\delta A/\delta g^{\mu\nu})\cdot B
+ A\cdot(\delta B/\delta g^{\mu\nu}) + \ldots$

For a \emph{nonlocal} action
\begin{equation}\label{eq:nonlocal-action-schematic}
  S = \int\sqrt{g}\,A\,F(\Box)\,B,
\end{equation}
the variation generates an additional term because $\Box$ depends on
$g^{\mu\nu}$:
\begin{equation}\label{eq:central-difficulty}
  \frac{\delta}{\delta g^{\mu\nu}}\big[A\,F(\Box)\,B\big]
  = \frac{\delta A}{\delta g^{\mu\nu}}\,F(\Box)\,B
  + A\,F(\Box)\,\frac{\delta B}{\delta g^{\mu\nu}}
  + A\,\frac{\delta F(\Box)}{\delta g^{\mu\nu}}\,B.
\end{equation}
The third term requires computing
$\delta F(\Box)/\delta g^{\mu\nu}$, which involves
$\delta\Box/\delta g^{\mu\nu}$.  Since $F$ is a function (not just
$\Box$ itself), this requires the ``$\alpha$-insertion formula''
from the Barvinsky--Vilkovisky formalism.

\subsection{Variation of $\sqrt{g}$}

\begin{proposition}[Fundamental metric variations]
\label{prop:metric-var}
\exactlit\;
\begin{align}
  \frac{\delta g_{\alpha\beta}}{\delta g^{\mu\nu}}
  &= -\frac{1}{2}(g_{\alpha\mu}g_{\beta\nu}
    + g_{\alpha\nu}g_{\beta\mu}),
  \label{eq:delta-g-lower}\\[4pt]
  \frac{\delta\sqrt{|g|}}{\delta g^{\mu\nu}}
  &= -\frac{1}{2}\sqrt{|g|}\,g_{\mu\nu},
  \label{eq:delta-sqrtg}\\[4pt]
  \frac{\delta R}{\delta g^{\mu\nu}}
  &= R_{\mu\nu} + g_{\mu\nu}\Box - \nabla_\mu\nabla_\nu,
  \label{eq:delta-R}\\[4pt]
  \frac{\delta R_{\alpha\beta}}{\delta g^{\mu\nu}}
  &= \frac{1}{2}\big(
    g_{\alpha\mu}\,\nabla_\nu\nabla_\beta
    + g_{\alpha\nu}\,\nabla_\mu\nabla_\beta
    + \nabla_\alpha\nabla_\mu\,g_{\beta\nu}
    + \nabla_\alpha\nabla_\nu\,g_{\beta\mu}
  \notag\\
  &\quad - \Box\,g_{\alpha(\mu}g_{\nu)\beta}
    - g_{\alpha\beta}\nabla_\mu\nabla_\nu
    - R_{\alpha(\mu\nu)\beta}
  \big).
  \label{eq:delta-Ricci}
\end{align}
These are standard identities; see
Parker--Toms~\cite{ParkerToms2009} App.~B.
\end{proposition}


\subsection{Variation of the covariant d'Alembertian}

\begin{proposition}[Variation of $\Box$ acting on a scalar $\phi$]
\label{prop:delta-box}
\fromlit{Barvinsky--Vilkovisky, Phys.\ Rep.\ \textbf{119} (1985);
Deser--Woodard 2007}
\begin{equation}\label{eq:delta-box}
  \delta[\Box\phi] = (\delta g^{\mu\nu})\nabla_\mu\nabla_\nu\phi
    + g^{\mu\nu}\,\delta\Gamma^\rho_{\mu\nu}\,\nabla_\rho\phi,
\end{equation}
where
\begin{equation}\label{eq:delta-Gamma}
  \delta\Gamma^\rho_{\mu\nu} = \frac{1}{2}g^{\rho\sigma}\big(
    \nabla_\mu\delta g_{\nu\sigma}
    + \nabla_\nu\delta g_{\mu\sigma}
    - \nabla_\sigma\delta g_{\mu\nu}\big).
\end{equation}
For the specific variation $\delta g^{\mu\nu} = h^{\mu\nu}$ and
working to all orders on curved background, the key identity is:
\begin{equation}\label{eq:delta-box-scalar}
  \frac{\delta(\Box\phi)}{\delta g^{\mu\nu}}
  = \nabla_\mu\nabla_\nu\phi
    + \frac{1}{2}g_{\mu\nu}\,g^{\rho\sigma}\nabla_\rho\nabla_\sigma\phi\cdot(\ldots)
    + \text{connection terms}.
\end{equation}
\end{proposition}

\begin{remark}[Operator-level variation]
\label{rem:op-level}
At the operator level (acting on an arbitrary tensor), the variation
$\delta\Box/\delta g^{\mu\nu}$ is itself a differential operator.
This is what makes the problem fundamentally different from varying
a local action: $\delta F(\Box)/\delta g^{\mu\nu}$ involves an
\emph{integral} over the spectrum of $\Box$.
\end{remark}


%%======================================================================
\section{The $\alpha$-Insertion Formula (Barvinsky--Vilkovisky)}
\label{sec:alpha-insertion}
%%======================================================================

This section documents the key technique for varying nonlocal actions.

\subsection{The fundamental identity}

\begin{theorem}[$\alpha$-insertion formula]
\label{thm:alpha-insertion}
\fromlit{Barvinsky--Vilkovisky, Nucl.\ Phys.\ B \textbf{282} (1987)
163; B \textbf{333} (1990) 471}

Let $F(\Box)$ be an analytic function of the covariant d'Alembertian
$\Box = g^{\mu\nu}\nabla_\mu\nabla_\nu$, and let $A$, $B$ be scalar
fields (or more generally, sections of some vector bundle).  Then the
variation of $A\,F(\Box)\,B$ with respect to $g^{\mu\nu}$ is:
\begin{equation}\label{eq:alpha-insertion}
  \boxed{
  \delta\big[A\,F(\Box)\,B\big]
  = (\delta A)\,F(\Box)\,B
  + A\,F(\Box)\,(\delta B)
  + \int_0^1\!\dd\alpha\;
    \big[e^{\alpha\Box}A\big]\,
    F'_\alpha(\Box)\,(\delta\Box)\,
    \big[e^{(1-\alpha)\Box}B\big],
  }
\end{equation}
where the notation $F'_\alpha$ denotes the ``$\alpha$-deformed
derivative'' defined via the integral representation:
\begin{equation}\label{eq:F-alpha-def}
  F'_\alpha(\Box) = \frac{F(\Box_1) - F(\Box_2)}{\Box_1 - \Box_2}
  \bigg|_{\substack{\Box_1 = \alpha\Box \\ \Box_2 = (1-\alpha)\Box}}.
\end{equation}
\end{theorem}

\begin{remark}[Simplification for entire functions]
\label{rem:entire-simplification}
If $F(z) = \sum_{n=0}^\infty c_n z^n$ is an entire function (as is the
case for the SCT form factors, by NT-2), then the $\alpha$-insertion
formula can be written as a sum over the Taylor coefficients:
\begin{equation}\label{eq:alpha-taylor}
  \int_0^1\dd\alpha\;\big[e^{\alpha\Box}A\big]\,
  (\delta\Box)\,\big[e^{(1-\alpha)\Box}B\big]
  \cdot F'_\alpha(\Box)
  = \sum_{n=1}^\infty c_n \sum_{k=0}^{n-1}
    (\Box^k A)\,(\delta\Box)\,(\Box^{n-1-k} B).
\end{equation}
This is the ``Leibniz rule for $F(\Box)$,'' generalizing the standard
Leibniz rule for $\Box^n$.
\end{remark}

\begin{remark}[Physical meaning]
\label{rem:alpha-meaning}
The parameter $\alpha \in [0,1]$ interpolates between ``$A$ evaluated
at the full nonlocal scale'' and ``$B$ evaluated at the full nonlocal
scale.''  The $\alpha$-insertion accounts for the fact that on a curved
background, the ``nonlocal smearing'' generated by $F(\Box)$ interacts
with the geometry at all intermediate scales.
\end{remark}

\gap\; The precise form of the $\alpha$-insertion formula for
\emph{tensor-valued} form factors (where $\Box$ acts on rank-2 tensors
rather than scalars) requires additional curvature correction terms.
On a general curved background, the commutator
$[\nabla_\mu, \nabla_\nu]V_\rho = R_{\mu\nu\rho}{}^\sigma V_\sigma$
generates Riemann tensor insertions.  These must be tracked in the
NT-4b derivation.


\subsection{Variation of $R\,F(\Box)\,R$}

\begin{proposition}[Variation of the $R^2$ nonlocal term]
\label{prop:var-RFR}
\fromlit{Deser--Woodard 2007; Nojiri--Odintsov 2008;
Calcagni--Modesto 2211.05606}

The variation of $S_{R^2} = \int\sqrt{g}\,R\,F(\Box)\,R$ with respect
to $g^{\mu\nu}$ gives:
\begin{equation}\label{eq:var-RFR}
\begin{split}
  \frac{1}{\sqrt{g}}\frac{\delta S_{R^2}}{\delta g^{\mu\nu}}
  &= R_{\mu\nu}\,F(\Box)\,R
    + R\,F(\Box)\,R_{\mu\nu}
    - \frac{1}{2}g_{\mu\nu}\,R\,F(\Box)\,R \\
  &\quad + 2\nabla_\mu\nabla_\nu[F(\Box)\,R]
    - 2g_{\mu\nu}\Box[F(\Box)\,R] \\
  &\quad + g_{\mu\nu}\nabla_\rho[R\,\nabla^\rho F(\Box)\,R]
    - \nabla_\mu[R\,\nabla_\nu F(\Box)\,R]
    - \nabla_\nu[R\,\nabla_\mu F(\Box)\,R] \\
  &\quad + \Theta_{\mu\nu}^{(R)}[R, F(\Box), R].
\end{split}
\end{equation}
The first line contains the ``naive'' variation of $A$ and $B$
plus the metric determinant.  The second line contains the
$\delta R/\delta g^{\mu\nu}$ contributions (generating $\nabla\nabla$
and $\Box$ terms).  The third and fourth lines contain the
$\alpha$-insertion terms $\Theta_{\mu\nu}^{(R)}$ from
$\delta\Box/\delta g^{\mu\nu}$:
\begin{equation}\label{eq:Theta-R}
  \Theta_{\mu\nu}^{(R)}
  = -\sum_{n=1}^\infty c_n \sum_{k=0}^{n-1}
    \Big[
      \nabla_\mu(\Box^k R)\,\nabla_\nu(\Box^{n-1-k} R)
      - \frac{1}{2}g_{\mu\nu}\,
        \nabla_\rho(\Box^k R)\,\nabla^\rho(\Box^{n-1-k} R)
    \Big],
\end{equation}
where $F(z) = \sum c_n z^n$.  In closed form (using the
$\alpha$-representation):
\begin{equation}\label{eq:Theta-R-closed}
  \Theta_{\mu\nu}^{(R)} = -\int_0^1\!\dd\alpha\;
    \Big[
      \nabla_\mu\Phi_\alpha\,\nabla_\nu\Psi_{1-\alpha}
      - \frac{1}{2}g_{\mu\nu}\,
        \nabla_\rho\Phi_\alpha\,\nabla^\rho\Psi_{1-\alpha}
    \Big],
\end{equation}
where $\Phi_\alpha = e^{\alpha\Box}R$ and
$\Psi_{1-\alpha} = F'_\alpha(\Box)\,e^{(1-\alpha)\Box}R$.
\end{proposition}

\begin{remark}[Local limit check]
When $F(\Box) \to \beta$ (constant), all the $\Theta$ terms vanish
(since $F' = 0$ for constant $F$), and we recover:
\begin{equation}
  \frac{1}{\sqrt{g}}\frac{\delta}{\delta g^{\mu\nu}}
  \int\sqrt{g}\,\beta R^2
  = 2\beta\,H_{\mu\nu},
\end{equation}
which matches~\eqref{eq:H-tensor}.  This is a mandatory consistency
check.
\end{remark}


\subsection{Variation of $\Weylsq F(\Box)\,\Weylsq$ (nonlocal Bach)}

The variation of the Weyl-squared term is more involved because the
Weyl tensor has four indices.

\begin{proposition}[Variation of the nonlocal Weyl term]
\label{prop:var-C2FC2}
\fromlit{Stelle 1978 (local limit);
Calcagni--Modesto 2211.05606 (nonlocal case);
Giacchini--de~Paula~Netto 2021}

Schematically, the variation of
$S_{C^2} = \int\sqrt{g}\,C_{\mu\nu\rho\sigma}\,F(\Box)\,C^{\mu\nu\rho\sigma}$
yields:
\begin{equation}\label{eq:var-C2}
\begin{split}
  \frac{1}{\sqrt{g}}\frac{\delta S_{C^2}}{\delta g^{\mu\nu}}
  &= 2\,B_{\mu\nu}^{(F)}
    - \frac{1}{2}g_{\mu\nu}\,C_{\alpha\beta\gamma\delta}\,
      F(\Box)\,C^{\alpha\beta\gamma\delta} \\
  &\quad + \Theta_{\mu\nu}^{(C)}[C_{\rho\sigma\alpha\beta},
    F(\Box), C^{\rho\sigma\alpha\beta}],
\end{split}
\end{equation}
where $B_{\mu\nu}^{(F)}$ is the \textbf{nonlocal Bach tensor}:
\begin{equation}\label{eq:nonlocal-bach}
  B_{\mu\nu}^{(F)} = \nabla^\rho\nabla^\sigma
    \big[F(\Box)\,C_{\mu\rho\nu\sigma}\big]
    + \frac{1}{2}R^{\rho\sigma}\,F(\Box)\,C_{\mu\rho\nu\sigma}
    + (\text{mixed terms}),
\end{equation}
generalizing the local Bach tensor~\eqref{eq:bach} by replacing
the curvature $C$ with $F(\Box)\,C$.

The $\Theta_{\mu\nu}^{(C)}$ terms are the $\alpha$-insertion
contributions from the variation of $\Box$ inside $F(\Box)$ when
it acts on the four-index Weyl tensor.  These are structurally
analogous to~\eqref{eq:Theta-R} but involve the full four-index
tensor structure.
\end{proposition}

\gap\; The explicit closed-form expression for
$\Theta_{\mu\nu}^{(C)}$ is not available in a single reference in the
$\{C^2, R^2\}$ basis.  Most references work in the
$\{R_{\mu\nu}^2, R^2\}$ or $\{G_{\mu\nu}, R^{\mu\nu}\}$ basis
(see Calcagni--Modesto~\cite{CalcagniModesto2022}), and the translation
to the Weyl basis requires using~\eqref{eq:Weyl-decomp} plus the
Gauss--Bonnet identity.  This translation must be performed carefully
in the NT-4b derivation.


%%======================================================================
\section{Nonlocal Gravity in the $G_{\mu\nu}\gamma(\Box)R^{\mu\nu}$ Basis}
\label{sec:NLG-basis}
%%======================================================================

An alternative approach, used extensively in the IDG literature, writes
the action in terms of the Einstein tensor.

\subsection{The simplified action}

\begin{convention}[Calcagni--Modesto simplified action]
\label{conv:CMN-action}
\fromlit{Calcagni--Modesto 2211.05606, Eq.~(simplifiedNLGActionF)}
With the choice $\gamma_2 = -2\gamma_0 \equiv \gamma(\Box)$ and
$\gamma_4 = 0$:
\begin{equation}\label{eq:NLG-simplified}
  S_{\mathrm{NLG}} = \frac{1}{2\kappa^2}\int\dd^4 x\,\sqrt{|g|}\,
  \big[R + G_{\mu\nu}\,\gamma(\Box)\,R^{\mu\nu}\big].
\end{equation}
\end{convention}

\begin{theorem}[Equations of motion in the $G\gamma R$ basis]
\label{thm:NLG-eom}
\fromlit{Calcagni--Modesto 2211.05606, Eq.~(EOMsNLG);
Calcagni--Modesto--Leszczyszyn 2018}
The field equations derived from~\eqref{eq:NLG-simplified} are:
\begin{equation}\label{eq:NLG-eom}
\begin{split}
  \kappa^2 T_{\mu\nu} = {}
  &e^{-r_*\Box}G_{\mu\nu}
    - \frac{1}{2}g_{\mu\nu}\,G_{\rho\sigma}\,\gamma(\Box)\,R^{\rho\sigma}
    + 2G^\rho{}_{(\mu}\,\gamma(\Box)\,G_{\nu)\rho} \\
  &+ g_{\mu\nu}\nabla^\rho\nabla^\sigma\gamma(\Box)\,G_{\rho\sigma}
    - 2\nabla^\rho\nabla_{(\mu}\gamma(\Box)\,G_{\nu)\rho} \\
  &+ \frac{1}{2}\big(
      G_{\mu\nu}\,\gamma(\Box)\,R
      + R\,\gamma(\Box)\,G_{\mu\nu}
    \big) \\
  &+ \Theta_{\mu\nu}(R_{\rho\sigma}, G^{\rho\sigma}),
\end{split}
\end{equation}
where $\Theta_{\mu\nu}$ contains the $\alpha$-insertion contributions.
For the Wataghin/minimal form factor $\gamma(\Box) = (e^{-r_*\Box}-1)/\Box$:
\begin{equation}\label{eq:Theta-CMN}
  \Theta_{\mu\nu}(R_{\rho\sigma}, G^{\rho\sigma})
  = -\int_0^{r_*}\!\dd q\;
    \tilde{\Theta}_{\mu\nu}\!\left[
      e^{-q\Box}R_{\rho\sigma},\;
      \frac{e^{-(r_*-q)\Box}-1}{\Box}\,G^{\rho\sigma}
    \right],
\end{equation}
with:
\begin{align}
  \tilde{\Theta}_{\mu\nu}^{\mathrm{sym}}(A_{\rho\sigma}, B^{\rho\sigma})
  &= -\nabla_\mu A_{\rho\sigma}\,\nabla_\nu B^{\rho\sigma}
    + \frac{1}{4}g_{\mu\nu}\,
      \nabla_\tau(A_{\rho\sigma}\nabla^\tau B^{\rho\sigma}
      + B^{\rho\sigma}\nabla^\tau A_{\rho\sigma}),
  \label{eq:Theta-sym}\\[6pt]
  \tilde{\Theta}_{\mu\nu}^{\mathrm{antisym}}(A_{\rho\sigma}, B^{\rho\sigma})
  &= \frac{1}{4}g_{\mu\nu}\,
      \nabla_\tau(A_{\rho\sigma}\nabla^\tau B^{\rho\sigma}
      - B^{\rho\sigma}\nabla^\tau A_{\rho\sigma})
  \notag\\
  &\quad + \nabla_\rho(A_{\mu\sigma}\nabla^\rho B^\sigma{}_\nu
    - B^\sigma{}_\nu\nabla^\rho A_{\mu\sigma})
  \notag\\
  &\quad + \nabla_\sigma(B_{\mu\rho}\nabla_\nu A^{\rho\sigma}
    - A^{\rho\sigma}\nabla_\nu B_{\mu\rho})
  \notag\\
  &\quad + \nabla_\rho(A_{\mu\sigma}\nabla_\nu B^{\sigma\rho}
    - B^{\sigma\rho}\nabla_\nu A_{\mu\sigma}).
  \label{eq:Theta-antisym}
\end{align}
Here indices $\mu$ and $\nu$ are implicitly symmetrized.
\end{theorem}

\begin{remark}[Recovery of Einstein's equations]
Setting $\gamma(\Box) = 0$ (i.e., $r_* = 0$), all nonlocal terms
vanish and~\eqref{eq:NLG-eom} reduces to $G_{\mu\nu} = \kappa^2 T_{\mu\nu}$.
\end{remark}

\begin{remark}[Ricci-flat solutions]
\fromlit{Calcagni--Modesto 2211.05606, \S~SecRicciFlat}
If $R_{\mu\nu} = 0$ (hence $G_{\mu\nu} = 0$), then all terms
in~\eqref{eq:NLG-eom} vanish and the vacuum Einstein solutions
(Schwarzschild, Kerr) remain exact solutions of the nonlocal theory
when $\gamma_4 = 0$.  This is an important structural property.
\end{remark}


\subsection{Mapping to the SCT $\{C^2, R^2\}$ basis}

\begin{proposition}[Basis translation]
\label{prop:basis-translation}
\needscheck\;
The SCT action~\eqref{eq:Squad} in the $\{C^2, R^2\}$ basis maps
to the general action~\eqref{eq:NLG-general} via:
\begin{align}
  \gamma_4 &= \frac{1}{16\pi^2}\,F_1\!\left(\frac{\Box}{\Lambda^2}\right),
  \label{eq:gamma4}\\
  \gamma_2 &= -\frac{2}{16\pi^2}\,F_1\!\left(\frac{\Box}{\Lambda^2}\right),
  \label{eq:gamma2}\\
  \gamma_0 &= \frac{1}{16\pi^2}\left[
    F_2\!\left(\frac{\Box}{\Lambda^2},\xi\right)
    + \frac{1}{3}F_1\!\left(\frac{\Box}{\Lambda^2}\right)
  \right].
  \label{eq:gamma0}
\end{align}
Note that the Calcagni--Modesto simplified form~\eqref{eq:NLG-simplified}
uses $\gamma_4 = 0$, which does \emph{not} match the SCT action (where
$\gamma_4 \neq 0$ corresponds to the Weyl term).  Therefore, the
equations~\eqref{eq:NLG-eom} cannot be used directly for the SCT case.
Instead, one must either:
\begin{enumerate}[nosep,label=(\alph*)]
  \item Use the general action~\eqref{eq:NLG-general} with all three
    form factors, or
  \item Work directly in the $\{C^2, R^2\}$ basis and vary each term
    separately.
\end{enumerate}
We recommend approach~(b) for NT-4b, as it maintains the clean
spin-2/spin-0 separation throughout.
\end{proposition}


%%======================================================================
\section{Key Identities for the Nonlinear Variation}
\label{sec:identities}
%%======================================================================

\subsection{Commutator of $\Box$ with curvature tensors}

\begin{proposition}[Commutation relations]
\label{prop:commutators}
\fromlit{Avramidi, ``Heat Kernel and Quantum Gravity'' (2000), Ch.~3;
Parker--Toms Ch.~6}
On curved backgrounds, the covariant d'Alembertian does not commute
with curvature tensors.  The key relations are:
\begin{align}
  [\nabla_\mu, \nabla_\nu]\phi &= 0
    \quad\text{(scalar)},
  \label{eq:comm-scalar}\\
  [\nabla_\mu, \nabla_\nu]V_\rho &= R_{\mu\nu\rho}{}^\sigma V_\sigma,
  \label{eq:comm-vector}\\
  [\nabla_\mu, \nabla_\nu]T_{\rho\sigma}
  &= R_{\mu\nu\rho}{}^\alpha T_{\alpha\sigma}
    + R_{\mu\nu\sigma}{}^\alpha T_{\rho\alpha},
  \label{eq:comm-tensor}\\
  [\Box, \nabla_\mu]\phi
  &= R_{\mu\nu}\nabla^\nu\phi,
  \label{eq:box-nabla-scalar}\\
  [\Box, \nabla_\mu]V_\nu
  &= R_{\mu\rho}\nabla^\rho V_\nu
    + R_{\mu\nu}{}^{\rho\sigma}\nabla_\rho V_\sigma
    + (\nabla^\rho R_{\mu\rho\nu}{}^\sigma) V_\sigma.
  \label{eq:box-nabla-vector}
\end{align}
\end{proposition}

\begin{remark}[Implications for NT-4b]
The commutator $[\Box, \nabla]$ generates Ricci and Riemann insertions
every time one moves $\Box$ past a covariant derivative.  In the
$\alpha$-insertion formula, this means the $\Theta_{\mu\nu}$ terms
receive curvature corrections at every order.  For the SCT form factors
(which are entire functions), these corrections are automatically
encoded in the integral representation~\eqref{eq:Theta-R-closed}.
\end{remark}


\subsection{The Lichnerowicz operator}

\begin{definition}[Lichnerowicz operator]
\label{def:lichnerowicz}
\fromlit{Calcagni--Modesto 2211.05606; Lichnerowicz 1961}
Acting on a symmetric rank-2 tensor $\psi_{\mu\nu}$:
\begin{equation}\label{eq:lichnerowicz}
  \Delta_L\psi_{\mu\nu}
  = 2R^\rho{}_{\mu\nu}{}^\tau\psi_{\tau\rho}
    + R_{\mu\rho}\psi^\rho{}_\nu
    + R_{\sigma\nu}\psi^\sigma{}_\mu
    - \Box\psi_{\mu\nu}.
\end{equation}
On a scalar: $\Delta_L\psi = -\Box\psi$.
\end{definition}

\begin{remark}[Role in stability]
\fromlit{Calcagni--Modesto 2211.05606, \S~stabi}
The Lichnerowicz operator encodes the curvature corrections to $\Box$
when acting on tensors.  In the nonlocal theory, replacing
$-\Box \to \Delta_L$ in the form factor gives a theory
$\mathcal{L} = R + G_{\mu\nu}\gamma(\Delta_L)R^{\mu\nu}$ that shares
the same linearized spectrum on Minkowski and inherits stability of
Ricci-flat solutions to all perturbative orders.
\end{remark}


\subsection{Integration by parts with $F(\Box)$}

\begin{proposition}[Self-adjointness of $\Box$]
\label{prop:self-adjoint}
\exactlit\;
On a closed manifold (or with vanishing boundary terms):
\begin{equation}\label{eq:box-selfadjoint}
  \int\sqrt{g}\,A\,\Box B\,\dd^4 x
  = \int\sqrt{g}\,(\Box A)\,B\,\dd^4 x
  = -\int\sqrt{g}\,\nabla_\mu A\,\nabla^\mu B\,\dd^4 x.
\end{equation}
More generally, for an entire function $F(\Box)$:
\begin{equation}\label{eq:F-selfadjoint}
  \int\sqrt{g}\,A\,F(\Box)\,B\,\dd^4 x
  = \int\sqrt{g}\,\big[F(\Box)\,A\big]\,B\,\dd^4 x.
\end{equation}
This self-adjointness is crucial for simplifying the field equations:
terms like $R_{\mu\nu}\,F(\Box)\,R + R\,F(\Box)\,R_{\mu\nu}$ can be
symmetrized.
\end{proposition}


%%======================================================================
\section{Consistency Checks for the Full Nonlinear Equations}
\label{sec:checks}
%%======================================================================

\subsection{Bianchi identity}

\begin{proposition}[Diffeomorphism invariance]
\label{prop:bianchi-full}
\exactlit\;
The full nonlinear equations must satisfy:
\begin{equation}\label{eq:bianchi-full}
  \nabla^\mu\left(\frac{1}{\sqrt{g}}\,
  \frac{\delta S}{\delta g^{\mu\nu}}\right) = 0,
\end{equation}
identically (off-shell), as a consequence of the diffeomorphism
invariance of $S$.  This is the fully nonlinear contracted Bianchi
identity.

In the local (Stelle) limit, this reduces to
$\nabla^\mu(G_{\mu\nu} + 2\alpha B_{\mu\nu} + 2\beta H_{\mu\nu}) = 0$,
which holds because $\nabla^\mu G_{\mu\nu} = 0$,
$\nabla^\mu B_{\mu\nu} = 0$, and $\nabla^\mu H_{\mu\nu} = 0$
individually.  For the nonlocal case, the
$\alpha$-insertion terms $\Theta_{\mu\nu}$ must also be covariantly
conserved:
\begin{equation}
  \nabla^\mu\Theta_{\mu\nu} = 0.
\end{equation}
\end{proposition}

\subsection{Trace equation}

\begin{proposition}[Conformal anomaly from the trace]
\label{prop:trace}
\fromlit{Deser--Duff--Isham 1976; Duff 1994}
The trace $g^{\mu\nu}(\delta S/\delta g^{\mu\nu})$ of the field
equations in the local limit gives:
\begin{equation}\label{eq:trace-local}
  g^{\mu\nu}(G_{\mu\nu} + 2\alpha B_{\mu\nu} + \beta H_{\mu\nu})
  = -R + 0 + \beta(-6\Box R) = -R - 6\beta\Box R.
\end{equation}
The Bach tensor contribution vanishes ($B_{\mu\nu}$ is traceless).
In the nonlocal case, the trace should reproduce the conformal anomaly
structure when evaluated on conformally flat backgrounds.
\end{proposition}


\subsection{Linearized limit}

\begin{proposition}[Recovery of NT-4a]
\label{prop:linearized-limit}
\oursign\;
When the full nonlinear equations are linearized around flat space
($g_{\mu\nu} = \delta_{\mu\nu} + h_{\mu\nu}$, $|h| \ll 1$), all
nonlinear curvature terms (including the $\Theta_{\mu\nu}$ contributions)
must vanish or reduce to their linearized forms, reproducing exactly:
\begin{align}
  \Pi_{\mathrm{TT}}(z) &= 1 + \frac{13}{60}\,z\,\hat{F}_1(z),\\
  \Pi_s(z,\xi) &= 1 + 6\left(\xi-\frac{1}{6}\right)^2 z\,\hat{F}_2(z,\xi).
\end{align}
This is the most important consistency check: NT-4b must contain
NT-4a as a special case.
\end{proposition}


\subsection{De Sitter solutions}

\begin{proposition}[Maximally symmetric solutions]
\label{prop:deSitter}
\fromlit{Calcagni--Modesto 2211.05606, \S~sectionRF;
Biswas--Mazumdar--Siegel 2006}
On a maximally symmetric background (de~Sitter or anti-de~Sitter):
\begin{equation}
  R_{\mu\nu\rho\sigma} = \frac{R_0}{12}
    (g_{\mu\rho}g_{\nu\sigma} - g_{\mu\sigma}g_{\nu\rho}),
  \quad
  R_{\mu\nu} = \frac{R_0}{4}g_{\mu\nu},
  \quad
  R = R_0 = \text{const}.
\end{equation}
The Weyl tensor vanishes: $C_{\mu\nu\rho\sigma} = 0$.  Therefore the
$\Weylsq F(\Box)\Weylsq$ term contributes nothing, and the field
equations reduce to:
\begin{equation}
  G_{\mu\nu} + \beta\,H_{\mu\nu}\big|_{\text{dS}}
  = \kappa^2 T_{\mu\nu},
\end{equation}
where $H_{\mu\nu}|_{\text{dS}}$ involves only constant-curvature
contributions.  Since $\Box R_0 = 0$ and $\nabla_\mu R_0 = 0$, the
$H_{\mu\nu}$ tensor simplifies to:
\begin{equation}
  H_{\mu\nu}\big|_{\text{dS}}
  = -\frac{1}{2}g_{\mu\nu}R_0^2 + 2R_0\cdot\frac{R_0}{4}g_{\mu\nu}
  = 0.
\end{equation}
Therefore, the de~Sitter condition becomes simply
$G_{\mu\nu} = \kappa^2 T_{\mu\nu}$ (or $G_{\mu\nu} + \Lambda g_{\mu\nu} = 0$
in vacuum).  De~Sitter spacetime with the correct cosmological constant
is an exact solution of the full nonlocal theory.
\end{proposition}


%%======================================================================
\section{Auxiliary Field Formulation}
\label{sec:auxiliary}
%%======================================================================

An alternative approach to deriving the field equations uses auxiliary
fields to localize the nonlocal action.

\begin{proposition}[Auxiliary field method]
\label{prop:auxiliary}
\fromlit{Calcagni--Modesto 2211.05606, Eq.~(AuxAction)}
Introduce a symmetric tensor $\phi_{\mu\nu}$ and write:
\begin{equation}\label{eq:aux-action}
  \tilde{S}[g,\phi] = \frac{1}{2\kappa^2}\int\dd^4 x\,\sqrt{|g|}\,
  \left[R + (2R_{\mu\nu} - \phi_{\mu\nu}
    + \frac{1}{2}g_{\mu\nu}\phi)\,\gamma(\Box)\,\phi^{\mu\nu}\right],
\end{equation}
where $\phi = g^{\mu\nu}\phi_{\mu\nu}$.  The $\phi_{\mu\nu}$
equation of motion yields $\phi_{\mu\nu} = G_{\mu\nu}$ on-shell,
recovering the original nonlocal action.

The advantage is that the $g_{\mu\nu}$ variation of
$\tilde{S}$ is \emph{algebraically simpler} because the nonlocal
operator $\gamma(\Box)$ now acts on the auxiliary field rather than
directly on curvature tensors.
\end{proposition}

\begin{remark}[Degrees of freedom]
\fromlit{Calcagni--Modesto 2211.05606}
The auxiliary field $\phi_{\mu\nu}$ is a symmetric rank-2 tensor
(10 components) constrained by the Bianchi identity (4 constraints)
to have 6 degrees of freedom.  Combined with the 2 degrees of
freedom of the graviton, the total is $2 + 6 = 8$ degrees of freedom,
matching Stelle's gravity counting.

In the nonlocal theory, the key difference is that the additional
6 degrees of freedom are \emph{not ghosts}: the propagator
$\propto \gamma\,e^{-H(\Box)}$ has no additional poles when $\gamma$
is an entire function of appropriate type.
\end{remark}


%%======================================================================
\section{Notation Table}
\label{sec:notation}
%%======================================================================

\begin{center}
\renewcommand{\arraystretch}{1.3}
\begin{longtable}{lll}
\toprule
\textbf{Symbol} & \textbf{Definition} & \textbf{Source} \\
\midrule
\endfirsthead
\toprule
\textbf{Symbol} & \textbf{Definition} & \textbf{Source} \\
\midrule
\endhead
$g_{\mu\nu}$ & spacetime metric, signature $(-,+,+,+)$ & Standard \\
$R^\rho{}_{\sigma\mu\nu}$ & Riemann tensor
  ($\partial_\mu\Gamma^\rho_{\nu\sigma} - \ldots$) & Standard \\
$R_{\mu\nu}$ & Ricci tensor ($R^\rho{}_{\mu\rho\nu}$) & Standard \\
$R$ & Ricci scalar ($g^{\mu\nu}R_{\mu\nu}$) & Standard \\
$G_{\mu\nu}$ & Einstein tensor ($R_{\mu\nu} - \frac{1}{2}gR$) & Standard \\
$C_{\mu\nu\rho\sigma}$ & Weyl tensor & Standard \\
$E_4$ & Gauss--Bonnet density & Eq.~\eqref{eq:GB} \\
$\Box$ & $g^{\mu\nu}\nabla_\mu\nabla_\nu$ (d'Alembertian) & Standard \\
$\Lambda$ & spectral action cutoff scale & SCT P3 \\
$\Lambda_{\mathrm{cc}}$ & cosmological constant & $S_{\mathrm{EH}}$ \\
$\kappa^2$ & $8\pi G$ (gravitational coupling) & Standard \\
$F_1(z)$ & Weyl form factor ($z = \Box/\Lambda^2$) & NT-1b \\
$F_2(z,\xi)$ & $R^2$ form factor & NT-1b \\
$\hat{F}_i(z)$ & normalized form factor ($\hat{F}_i(0)=1$) & NT-4a \\
$\alpha_C$ & $16\pi^2 F_1(0) = 13/120$ & Phase~3 \\
$\alpha_R(\xi)$ & $16\pi^2 F_2(0) = 2(\xi-1/6)^2$ & Phase~3 \\
$c_1, c_2$ & Stelle basis coefficients & Phase~3 \\
$\varphi(z)$ & master function
  $e^{-z/4}\sqrt{\pi/z}\,\erfi(\sqrt{z}/2)$ & NT-1 \\
$\xi$ & Higgs non-minimal coupling & SM \\
$B_{\mu\nu}$ & Bach tensor (local) & Eq.~\eqref{eq:bach} \\
$B_{\mu\nu}^{(F)}$ & nonlocal Bach tensor & Eq.~\eqref{eq:nonlocal-bach} \\
$H_{\mu\nu}$ & scalar curvature tensor (local) & Eq.~\eqref{eq:H-tensor} \\
$\Theta_{\mu\nu}$ & $\alpha$-insertion contribution & BV formalism \\
$\gamma_0, \gamma_2, \gamma_4$ & form factors in Riemann basis & CMN \\
$\Delta_L$ & Lichnerowicz operator & Eq.~\eqref{eq:lichnerowicz} \\
$\Pi_{\mathrm{TT}}(z)$ & spin-2 propagator denominator & NT-4a \\
$\Pi_s(z,\xi)$ & spin-0 propagator denominator & NT-4a \\
$m_2$ & $\Lambda\sqrt{60/13}$ (spin-2 mass) & NT-4a \\
$m_0$ & $\Lambda/\sqrt{6(\xi-1/6)^2}$ (scalar mass) & NT-4a \\
$N_s, N_f, N_v$ & SM multiplicities ($4, 45, 12$) & CPR \\
$N_D$ & Dirac fermion count ($N_f/2 = 22.5$) & Phase~3 \\
\bottomrule
\end{longtable}
\end{center}


%%======================================================================
\section{Gap Analysis}
\label{sec:gaps}
%%======================================================================

\begin{enumerate}[label=\textbf{G\arabic*.}]
\item \textbf{Explicit $\Theta_{\mu\nu}^{(C)}$ in the Weyl basis.}
  The $\alpha$-insertion terms for the Weyl-squared variation
  $\delta(\Weylsq F(\Box)\Weylsq)/\delta g^{\mu\nu}$ are not
  available in closed form in a single published reference in the
  $\{C^2, R^2\}$ basis.  The Calcagni--Modesto result is in the
  $\{G_{\mu\nu}, R^{\mu\nu}\}$ basis.  Translation requires
  careful application of the Weyl decomposition and Gauss--Bonnet
  identity.

\item \textbf{Tensor-valued $\alpha$-insertion.}
  The $\alpha$-insertion formula (Theorem~\ref{thm:alpha-insertion}) is
  stated for scalar-valued quantities $A$, $B$.  When $\Box$ acts on
  rank-4 tensors (as in $F(\Box)\,C_{\mu\nu\rho\sigma}$), the
  commutator $[\Box, \nabla]$ generates Riemann tensor insertions at
  every order.  The systematic treatment requires the full
  Barvinsky--Vilkovisky covariant perturbation theory to second order
  in curvature.

\item \textbf{Non-commutativity of $F_1$ and $F_2$ form factors.}
  In the SCT action, $F_1$ and $F_2$ are different functions of
  $\Box/\Lambda^2$.  When varying $g^{\mu\nu}$, the variations of
  $F_1(\Box)$ and $F_2(\Box)$ generate different $\Theta$-type terms.
  These must be computed separately and summed.

\item \textbf{Boundary terms.}
  The self-adjointness identity~\eqref{eq:F-selfadjoint} assumes
  vanishing boundary contributions.  For the FLRW reduction (NT-4c),
  boundary terms at the initial singularity or at spatial infinity
  may need careful treatment.

\item \textbf{Sign convention propagation.}
  The Stelle equations~\eqref{eq:stelle-eom} use the convention
  $G_{\mu\nu} + 2\alpha B_{\mu\nu} + \beta H_{\mu\nu} = \kappa^2 T_{\mu\nu}$.
  The mapping $\alpha \to F_1(0)/(16\pi^2)$,
  $\beta \to F_2(0)/(16\pi^2)$ must include all factors of 2 from the
  Weyl decomposition.  This is the most common source of sign errors
  in the literature.

\item \textbf{The Gauss--Bonnet term with nonlocal form factor.}
  If one works in the full $\{R_{\mu\nu\rho\sigma}^2, \Ricsq, R^2\}$
  basis, the Gauss--Bonnet topological identity
  $\delta(\sqrt{g}\,E_4)/\delta g^{\mu\nu} = 0$ holds only when the
  form factor is unity (local case).  For
  $\int\sqrt{g}\,R_{\mu\nu\rho\sigma}\,\gamma_4(\Box)\,R^{\mu\nu\rho\sigma}$
  with $\gamma_4 \neq \text{const}$, the Gauss--Bonnet cancellation
  does \emph{not} hold because $[\gamma_4(\Box), \delta/\delta g^{\mu\nu}]
  \neq 0$.  This means the basis reduction must be done \emph{before}
  applying the form factor, not after.

\item \textbf{Connection to the spectral action.}
  The form factors $F_1$ and $F_2$ are derived from the heat kernel
  expansion of $\Tr(f(D^2/\Lambda^2))$.  The nonlocal effective action
  is obtained by resumming the heat kernel to second order in curvature
  (Barvinsky--Vilkovisky CPT-II).  The field equations are the
  Euler--Lagrange equations of this resummed action.  The consistency
  between the heat-kernel resummation and the variational equations
  must be verified.
\end{enumerate}


%%======================================================================
\section{Strategy for the NT-4b Derivation}
\label{sec:strategy}
%%======================================================================

Based on the literature survey, we recommend the following strategy
for the NT-4b derivation:

\begin{enumerate}[label=\textbf{Step \arabic*.}]
\item \textbf{Work in the $\{C^2, R^2\}$ basis.}
  This maintains the clean spin-2/spin-0 separation throughout and
  allows direct comparison with the NT-4a linearized results.

\item \textbf{Vary $S_{R^2}$ and $S_{C^2}$ separately.}
  Use the formulas from \S\ref{sec:nonlocal-variation} for each term.
  The $R^2$ variation is cleaner (scalar $\times$ scalar) and should
  be done first as a warm-up.

\item \textbf{Use the $\alpha$-insertion formula for $\Theta$ terms.}
  For each form factor, compute the $\alpha$-insertion contribution
  using the Taylor expansion~\eqref{eq:alpha-taylor} (which is valid
  because the SCT form factors are entire by NT-2).

\item \textbf{Verify the local limit.}
  Set $F_i \to F_i(0)$ and confirm recovery of the Stelle
  equations~\eqref{eq:stelle-eom} with the correct coefficients
  $\alpha = \alpha_C/(16\pi^2)$, $\beta = \alpha_R/(16\pi^2)$.

\item \textbf{Verify the Bianchi identity.}
  Check $\nabla^\mu(\delta S/\delta g^{\mu\nu}) = 0$ at the
  full nonlinear level.

\item \textbf{Linearize and recover NT-4a.}
  Expand around flat space and verify recovery of
  $\Pi_{\mathrm{TT}}$ and $\Pi_s$.

\item \textbf{Check de~Sitter and Ricci-flat solutions.}
  Verify that de~Sitter spacetime and Schwarzschild remain solutions.

\item \textbf{Write the final field equations in canonical form.}
  Present the result as:
  \begin{equation}
    G_{\mu\nu} + 2\alpha\,B_{\mu\nu}^{(F_1)}
    + \beta\,H_{\mu\nu}^{(F_2)}
    + \Theta_{\mu\nu}^{(1)} + \Theta_{\mu\nu}^{(2)}
    = \kappa^2 T_{\mu\nu}.
  \end{equation}
\end{enumerate}


%%======================================================================
\begin{thebibliography}{99}
%%======================================================================

\bibitem{Stelle1977}
  K.~S.~Stelle,
  ``Renormalization of higher-derivative quantum gravity,''
  Phys.\ Rev.\ D \textbf{16} (1977) 953.

\bibitem{Stelle1978}
  K.~S.~Stelle,
  ``Classical gravity with higher derivatives,''
  Gen.\ Rel.\ Grav.\ \textbf{9} (1978) 353--371.

\bibitem{BV1985}
  A.~O.~Barvinsky, G.~A.~Vilkovisky,
  ``The generalized Schwinger-DeWitt technique in gauge theories and
  quantum gravity,''
  Phys.\ Rep.\ \textbf{119} (1985) 1--74.

\bibitem{BV1987}
  A.~O.~Barvinsky, G.~A.~Vilkovisky,
  ``Beyond the Schwinger-DeWitt technique: Converting loops into
  trees and in-in currents,''
  Nucl.\ Phys.\ B \textbf{282} (1987) 163--188.

\bibitem{BV1990}
  A.~O.~Barvinsky, G.~A.~Vilkovisky,
  ``Covariant perturbation theory (II): Second order in the
  curvature. General algorithms,''
  Nucl.\ Phys.\ B \textbf{333} (1990) 471--511.

\bibitem{BV1990b}
  A.~O.~Barvinsky, G.~A.~Vilkovisky,
  ``Covariant perturbation theory (III): Spectral representations
  of the third-order form factors,''
  Nucl.\ Phys.\ B \textbf{333} (1990) 512--524.

\bibitem{CZ2013}
  A.~Codello, O.~Zanusso,
  ``On the non-local heat kernel expansion,''
  J.\ Math.\ Phys.\ \textbf{54} (2013) 013513,
  arXiv:1203.2034 [hep-th].

\bibitem{CPR2009}
  A.~Codello, R.~Percacci, C.~Rahmede,
  ``Investigating the ultraviolet properties of gravity with a
  Wilsonian renormalization group equation,''
  Ann.\ Phys.\ \textbf{324} (2009) 414,
  arXiv:0805.2909 [hep-th].

\bibitem{BMS2006}
  T.~Biswas, A.~Mazumdar, W.~Siegel,
  ``Bouncing universes in string-inspired gravity,''
  JCAP \textbf{0603} (2006) 009,
  hep-th/0508194.

\bibitem{Modesto2012}
  L.~Modesto,
  ``Super-renormalizable quantum gravity,''
  Phys.\ Rev.\ D \textbf{86} (2012) 044005,
  arXiv:1107.2403 [hep-th].

\bibitem{CalcagniModesto2022}
  G.~Calcagni, L.~Modesto,
  ``Classical and quantum nonlocal gravity,''
  arXiv:2211.05606 [hep-th] (2022).

\bibitem{CalcagniModesto2018}
  G.~Calcagni, L.~Modesto, G.~Nardelli,
  ``Initial conditions and degrees of freedom of non-local gravity,''
  JHEP \textbf{05} (2018) 087,
  arXiv:1803.00561 [hep-th].

\bibitem{CalcagniModesto2018b}
  G.~Calcagni, L.~Modesto, G.~Nardelli,
  ``Non-perturbative spectrum of non-local gravity,''
  arXiv:1803.07848 [hep-th].

\bibitem{Buoninfante2020}
  L.~Buoninfante, G.~Lambiase, Y.~Miyashita, W.~Takebe,
  M.~Yamaguchi,
  ``Generalized ghost-free propagators in nonlocal field theories,''
  Phys.\ Rev.\ D \textbf{101} (2020) 084019,
  arXiv:2001.07830 [hep-th].

\bibitem{GiacchiniNetto2021}
  B.~L.~Giacchini, T.~de~Paula~Netto,
  ``Higher-order regularity in local and nonlocal quantum gravity,''
  Eur.\ Phys.\ J.\ C \textbf{81} (2021) 741.

\bibitem{GiacchiniNetto2018}
  B.~L.~Giacchini, T.~de~Paula~Netto,
  ``Effective delta sources and regularity in
  higher-derivative and ghost-free gravity,''
  JCAP \textbf{07} (2019) 013,
  arXiv:1806.05664 [gr-qc].

\bibitem{Tomboulis1997}
  E.~T.~Tomboulis,
  ``Super-renormalizable gauge and gravitational theories,''
  hep-th/9702080 (1997).

\bibitem{Avramidi2000}
  I.~G.~Avramidi,
  ``Heat Kernel and Quantum Gravity,''
  Springer, Lecture Notes in Physics \textbf{M64} (2000).

\bibitem{ParkerToms2009}
  L.~Parker, D.~Toms,
  ``Quantum Field Theory in Curved Spacetime: Quantized Fields
  and Gravity,''
  Cambridge University Press (2009).

\bibitem{Briscese2019}
  F.~Briscese, G.~Calcagni, L.~Modesto,
  ``Nonlinear stability in nonlocal gravity,''
  Phys.\ Rev.\ D \textbf{99} (2019) 084041,
  arXiv:1901.03267 [gr-qc].

\bibitem{DeserWoodard2007}
  S.~Deser, R.~P.~Woodard,
  ``Nonlocal cosmology,''
  Phys.\ Rev.\ Lett.\ \textbf{99} (2007) 111301.

\bibitem{CC1997}
  A.~H.~Chamseddine, A.~Connes,
  ``The spectral action principle,''
  Commun.\ Math.\ Phys.\ \textbf{186} (1997) 731,
  hep-th/9606001.

\bibitem{NT4aLit}
  D.~Alfyorov,
  ``NT-4a: Literature extraction \& convention fixing for linearized
  nonlocal field equations,''
  SCT Theory internal document (2026).

\bibitem{BGM2012}
  T.~Biswas, E.~Gerber, A.~Mazumdar,
  ``Towards singularity- and ghost-free theories of gravity,''
  Phys.\ Rev.\ Lett.\ \textbf{108} (2012) 031101,
  arXiv:1112.5430 [gr-qc].

\end{thebibliography}

\end{document}
